% from ega4-11-fr.tex
\oldpage[IV-3]{160}

point $x'\in X'$ dont les projections dans $X_1$ et $Y'$ sont respectivement
$x_1$ et $y'$ (\textbf I, 3.4.7)~; comme par hypothèse $\mathcal F'$ est
$f'$-plat au point $x'$, et que $\mathcal O_{y_1}$ est intégralement clos, on
peut appliquer (11.6.1), et on en déduit que $\mathcal F_1$ est $f_1$-plat au
point $x_1$, d'où le théorème dans ce cas.

\medskip
\noindent\textbf{(ii)} \emph{Cas où $A$ est réduit et noethérien.} --- Soient
$\mathfrak p_i$ $(1\leq i\leq m)$ les idéaux minimaux de $A$~; comme $A$ est
réduit, il s'identifie canoniquement à un sous-anneau du produit des
$A_i=A/\mathfrak p_i$, qui sont des anneaux locaux noethériens~; posant
$Y_i=\operatorname{Spec}(A_i)$, $X_i=X\times_Y Y_i$, $f_i=f_{(Y_i)}$, il
résulte de (11.5.2) qu'il suffit de démontrer que pour chaque $i$,
$\mathcal F_i=\mathcal F\otimes_Y Y_i$ est $f_i$-plat en tout point $x_i$ de
$X_i$ dont les projections dans $X$ et $Y_i$ sont $x$ et le point fermé $y_i$
de $Y_i$ respectivement. Or, comme $A_i$ est un anneau local noethérien intègre,
il existe dans son corps des fractions $K_i$ un sous-anneau $A''_i$ contenant
$A_i$, qui soit une $A_i$-algèbre \emph{finie} (donc un anneau semi-local
noethérien) et dont les anneaux locaux soient \emph{géométriquement
unibranches} ((6.15.5) et $(\mathbf 0,23.2.5)$). Comme les idéaux maximaux de
$A''_i$ sont alors nécessairement au-dessus de l'idéal maximal de $A_i$, on
déduit encore de (\textbf I, 3.4.7) et de (11.5.2) qu'il suffit, en posant
$Y''_i=\operatorname{Spec}(A''_i)$, $X''_i=X\times_Y Y''_i$,
$f''_i=f_{(Y''_i)}$, de prouver que
$\mathcal F''_i=\mathcal F\otimes_Y Y''_i$ est $f''_i$-plat en tout point
$x''_i$ de $X''_i$ dont les projections dans $X$ et $Y''_i$ sont $x$ et le
point fermé $y''_i$ de $Y''_i$ respectivement. Or, soit $A'$ un anneau de
valuation discrète pour $K_i$ dominant $A''_i$, et soit $x'$ un point de $X'$
dont les projections dans $X''_i$ et dans $Y'$ sont $x''_i$ et $y'$
respectivement (\textbf I, 3.4.7)~; comme $\mathcal O_{y''_i}$ est
géométriquement unibranche, on peut encore appliquer (11.6.1) et on en déduit
bien que $\mathcal F''_i$ est $f''_i$-plat au point $x''_i$.
\end{proof}

\begin{remarks}[11.8.2]
\label{IV.11.8.2-fr}
\begin{enumerate}[label=(\roman*)]
  \item Dans l'énoncé de (11.8.1), on peut se borner à supposer que la condition
  sur $\mathcal F'$ est vérifiée pour les anneaux de valuation $A'$ \emph{complets}
  et dont le corps résiduel est \emph{algébriquement clos}. On sait en effet que
  tout anneau de valuation $A'$ est dominé par un tel anneau $A''$
  (\textbf{II}, 7.1.2), et que si $A'$ est un anneau de valuation discrète, on
  peut supposer qu'il en est de même de $A''$
  $(\mathbf 0_{\mathrm{III}},10.3.1)$.
  \item La démonstration de (11.8.1) se simplifie lorsqu'on suppose non seulement
  que $A$ est intègre et noethérien, mais que son complété $\widehat A$ est lui
  aussi \emph{intègre}. Remplaçant $X$ par $X\otimes_A\widehat A$ et raisonnant
  comme dans la démonstration de (11.5.3), on peut en effet se ramener alors à
  prouver (11.8.1) lorsque $A=\mathcal O_y$ est intègre, noethérien et
  \emph{complet}. Or, on sait (\textbf{II}, 7.1.7) qu'un tel anneau $A$ est
  dominé par un anneau de valuation discrète complet~; la conclusion résulte
  donc directement de (11.5.2).
\end{enumerate}
\end{remarks}

\subsection{Familles séparantes et universellement séparantes d'homomorphismes
de faisceaux de modules}
\label{subsection:IV.11.9-fr}

\begin{env}[11.9.1]
\label{IV.11.9.1-fr}
Soient $X$ un préschéma, $(f_{\lambda})_{\lambda\in L}$ une famille de morphismes
$f_{\lambda}:Z_{\lambda}\to X$, $\mathcal F$ un $\mathcal O_X$-Module
quasi-cohérent~; pour tout $\lambda\in L$, supposons donnés un
$\mathcal O_{Z_{\lambda}}$-Module quasi-cohérent $\mathcal G_{\lambda}$ et un
homomorphisme
\begin{equation}
u_{\lambda}:\mathcal F\longrightarrow(f_{\lambda})_*(\mathcal G_{\lambda}).
\label{IV.11.9.1.1-fr}
\tag{11.9.1.1}
\end{equation}
\end{env}
% IV3_B23_P160_END

